\section{Conclusion}


The sovereign key architecture on Lux Network ensures that service providers
cannot access user data.  Users derive root keys from passphrases via
Argon2id, encrypt data with ML-KEM-1024 + AES-256-GCM envelopes, and
sign with ML-DSA-87---all at NIST post-quantum Security Level 5.
Threshold MPC (CGGMP21/FROST) distributes signing authority among
user-controlled parties with zero operator shards.  FHE enables
computation on encrypted data without decryption.  The architecture
uses the same cryptographic primitives as the regulated custody model
but inverts the trust hierarchy: the user is the root of trust, not
the operator.

% ======================================================================
\begin{thebibliography}{99}

\bibitem{fips203}
  National Institute of Standards and Technology,
  ``FIPS 203: Module-Lattice-Based Key-Encapsulation Mechanism Standard (ML-KEM),''
  August 2024.

\bibitem{fips204}
  National Institute of Standards and Technology,
  ``FIPS 204: Module-Lattice-Based Digital Signature Standard (ML-DSA),''
  August 2024.

\bibitem{fips197}
  National Institute of Standards and Technology,
  ``FIPS 197: Advanced Encryption Standard (AES),''
  November 2001.

\bibitem{argon2}
  A.~Biryukov, D.~Dinu, D.~Khovratovich,
  ``Argon2: New Generation of Memory-Hard Functions for Password Hashing and Other Applications,''
  RFC 9106, September 2021.

\bibitem{cggmp21}
  R.~Canetti, R.~Gennaro, S.~Goldfeder, N.~Makriyannis, U.~Peled,
  ``UC Non-Interactive, Proactive, Threshold ECDSA with Identifiable Aborts,''
  in \emph{Proceedings of the 2020 ACM CCS}, 2020.

\bibitem{frost}
  C.~Komlo and I.~Goldberg,
  ``FROST: Flexible Round-Optimized Schnorr Threshold Signatures,''
  in \emph{Selected Areas in Cryptography (SAC)}, 2020.

\bibitem{lux-messaging}
  Lux Industries,
  ``Secure Messaging Protocol for Lux Network,''
  Technical Report, 2026.

\end{thebibliography}

